\chapter{Project Overview}
\label{ch:project-overview}

\section{Introduction}

\textbf{Expo Snips} is a developer tooling product that generates complete
React Native / Expo screens from a short identifier. It ships as two independent
developer surfaces that share a single generation engine:

\begin{itemize}
  \item a command-line interface published as \texttt{@expo-snips/cli};
  \item a Visual Studio Code extension published as \texttt{expo-snippets}
        (display name ``Expo Snips'').
\end{itemize}

Both surfaces render versioned ``screen recipes'' from a registry into the
developer's existing project, then optionally install the recipe's npm
dependencies and record the installation in a project-level manifest file.

\textbf{Status:} IMPLEMENTED --- see \texttt{packages/cli/package.json},
\texttt{packages/vscode-extension/package.json} and \texttt{packages/core/src}.

\section{Repository at a glance}

\begin{table}[H]
\centering
\caption{Repository identity and metrics (all values verified from the working tree)}
\begin{tabular}{@{}L{5.4cm}L{9cm}@{}}
\toprule
\textbf{Attribute} & \textbf{Value} \\
\midrule
Root package name & \texttt{expo-snips} (private workspace root) \\
Workspace root version & 0.1.0 \\
Package manager & pnpm 11.9 (workspace declared in \texttt{pnpm-workspace.yaml}) \\
Workspace packages & 4 (\texttt{packages/*}) \\
Languages & TypeScript (all packages), JavaScript (build scripts, \texttt{.mjs}) \\
Node.js requirement & $\geq$ 18 (\texttt{engines.node}) \\
TypeScript version & 5.5 \\
Screen recipes & 24 (\texttt{packages/registry/data}) \\
Handlebars templates & 48 (\texttt{packages/registry/templates}) \\
Automated test files & 0 \\
CI workflows & 2 (\texttt{deploy-registry.yml}, \texttt{publish.yml}) \\
Root README & none \\
Default branch & \texttt{master} (plus \texttt{origin/gh-pages}) \\
Commit count at analysis time & 10 \\
\bottomrule
\end{tabular}
\end{table}

\section{Product category}

\textbf{Development tooling / code generation.} Expo Snips is not a SaaS
application, not a hosted platform and not a component library. It is a local
developer tool that writes files into the user's own repository.

\textbf{Status:} IMPLEMENTED --- the CLI writes to \texttt{process.cwd()} and the
extension writes to the VS Code workspace root; there is no server component
(\texttt{packages/cli/src/utils/writeFiles.ts},
\texttt{packages/vscode-extension/src/generator.ts}).

\section{Problems addressed}

The repository's own documentation (\texttt{packages/vscode-extension/README.md},
lines 27--39) states the problem: building Expo / React Native apps requires
re-writing the same screens repeatedly --- authentication flows, app chrome,
e-commerce views, onboarding and empty/error states --- each requiring multiple
files, careful styling and several library integrations.

The templates in the repository confirm this scope: recipes generate a screen
file, form/child components and hooks, with validation, loading states, error
handling and keyboard/safe-area handling
(\texttt{packages/registry/templates/auth/login/*}).

\section{Solution}

A versioned catalog of screen recipes, rendered by a shared Handlebars engine
that adapts output to the target project's detected stack (styling, router,
UI primitives, icon library, TypeScript or JavaScript).

\textbf{Status:} IMPLEMENTED --- \texttt{packages/core/src/generate.ts}
(\texttt{getFilesForRecipe}, conditional \texttt{when} filtering) and
\texttt{packages/cli/src/utils/detectProject.ts}.

\section{Primary users}

\begin{itemize}
  \item Expo / React Native application developers building mobile apps.
  \item Developers who prefer a terminal workflow (CLI) or an editor-native
        workflow (VS Code completion, command palette, status bar).
  \item Teams that want a consistent, version-tracked catalog of screen
        implementations (via the manifest and the \texttt{update} command).
\end{itemize}

\textbf{Status:} IMPLEMENTED (developer as the only actor) /
INFERRED (team and agency usage, which the repository supports technically but
does not demonstrate).

\section{High-level capabilities}

\begin{longtable}{@{}L{4.6cm}L{9.8cm}@{}}
\caption{Capability summary}\\
\toprule
\textbf{Capability} & \textbf{Description} \\
\midrule
\endfirsthead
\toprule
\textbf{Capability} & \textbf{Description} \\
\midrule
\endhead
Screen generation & Render a recipe into one or more project files. \\
Stack adaptation & Detect router, styling, UI, icons and TS/JS; render conditionally. \\
Interactive selection & Prompt for unresolved options (CLI) or fall back to settings (VS Code). \\
Dependency installation & Run the project's package manager to add recipe dependencies. \\
Primitive resolution & Recursively add shadcn-style UI primitives when \texttt{ui=shadcn}. \\
Manifest tracking & Record installs in \texttt{.expo-snips.json}. \\
Versioned updates & Re-render tracked screens and preview a line diff before writing. \\
Catalog distribution & Build and serve the registry as static JSON; bundle the catalog offline. \\
Editor integration & Inline \texttt{snip:} completion, command palette entry and status bar item. \\
Project bootstrap & Create a new Expo project through \texttt{create-expo-app}. \\
\bottomrule
\end{longtable}

\section{What this product is not}

To prevent misreading of later sections, the following are explicitly \textbf{not}
present in the repository:

\begin{itemize}
  \item No backend service and no REST/GraphQL API server
        (\textbf{NOT APPLICABLE}).
  \item No database, ORM, migrations or schema files
        (\textbf{NOT APPLICABLE}). The only persisted state is the
        \texttt{.expo-snips.json} manifest.
  \item No user accounts, authentication, authorization or sessions in the product
        itself (\textbf{NOT APPLICABLE}). Authentication \emph{templates} are
        generated for the user's app; see Chapter~\ref{ch:authentication}.
  \item No web application, marketing site or dashboard.
  \item No tests, linters or formatter configuration.
  \item No Docker, Terraform or other infrastructure-as-code.
\end{itemize}

\section{Navigating this document}

Chapter~\ref{ch:executive-summary} is written for non-technical stakeholders.
Chapters~\ref{ch:brd}--\ref{ch:prd} cover business and product requirements.
Chapters~\ref{ch:system-architecture}--\ref{ch:frontend-architecture} cover
architecture and design. Chapter~\ref{ch:cli-reference} is the operational
command reference, Chapter~\ref{ch:api-documentation} documents every callable
interface, and Chapters~\ref{ch:known-issues}--\ref{ch:technical-debt} record
verified problems and debt.
